\documentclass[10pt]{article}
\usepackage[T1]{fontenc}
\usepackage[utf8]{inputenc}
\usepackage[french]{babel}
\usepackage{amsmath,amssymb,mathrsfs,mathpazo}
\usepackage{geometry}
\geometry{a4paper,margin=2cm}
\usepackage{graphicx,array,tabularx,tabularray,colortbl,multirow,multicol,booktabs,enumitem,pifont,ulem,fancybox,diagbox,cancel}
\usepackage{pstricks,pstricks-add,pst-plot,pst-tree,pst-node,pst-text,pst-func,pst-math,pst-bspline,pst-3dplot,pst-eucl,pst-solides3d}
\let\pscalculate\undefined
\usepackage{pst-bezier}
\usepackage{tikz,pgfplots,tkz-tab}
\usetikzlibrary{arrows,arrows.meta,patterns,calc,positioning,trees,decorations.markings,shapes.geometric,shapes.misc,shapes.arrows,intersections,decorations.pathmorphing,decorations.pathreplacing,backgrounds,shadows,babel}
\pgfplotsset{compat=1.18}
\usepgfplotslibrary{fillbetween}
\usepackage[np]{numprint}
\usepackage{listings,makecell,pas-tableur,twoopt}
\let\barre\undefined
\usepackage[pstricks]{tablvar}
\definecolor{darkgreen}{RGB}{0,100,0}
\providecommand{\red}{\color{red}}
\providecommand{\blue}{\color{blue}}
\providecommand{\green}{\color{green}}
\providecommand{\black}{\color{black}}
\providecommand{\white}{\color{white}}
\providecommand{\gray}{\color{gray}}
\providecommand{\cyan}{\color{cyan}}
\providecommand{\magenta}{\color{magenta}}
\providecommand{\yellow}{\color{yellow}}
\providecommand{\violet}{\color{violet}}
\providecommand{\orange}{\color{orange}}
\providecommand{\euro}{\texteuro}
\providecommand{\EUR}{\texteuro}
\providecommand{\ieme}{\textsuperscript{e}}
\providecommand{\degres}{\ensuremath{^\circ}}
\providecommand{\degree}{\ensuremath{^\circ}}
\providecommand{\vv}[1]{\overrightarrow{#1}}
\providecommand{\Coord}[1]{#1}
\providecommand{\decosix}{}
\providecommand{\starredbullet}{\ensuremath{\bullet}}
\pagestyle{empty}
\setlength{\parindent}{0pt}
\usepackage{adjustbox}\DeclareUnicodeCharacter{2236}{:}\DeclareUnicodeCharacter{2190}{\ensuremath{\leftarrow}}
\providecommand{\R}{\mathbb{R}}
\providecommand{\N}{\mathbb{N}}
\providecommand{\Z}{\mathbb{Z}}
\providecommand{\Q}{\mathbb{Q}}
\providecommand{\D}{\mathbb{D}}
\providecommand{\C}{\mathbb{C}}
\providecommand{\e}{{\mathrm{e}}}
\providecommand{\vect}[1]{\overrightarrow{#1}}
\providecommand{\vectt}[1]{\overrightarrow{\mathrm{#1}}}
\providecommand{\vvt}[1]{\overrightarrow{\mathrm{#1}}}
\providecommand{\np}[2][]{{#2~#1}}
\providecommand{\vv}[1]{{\overrightarrow{#1}}}
\renewcommand{\d}{{\mathrm{d}}}
\providecommand{\degres}{{^\circ}}
\providecommand{\degre}{{^\circ}}
\providecommand{\blue}{\color{blue}}
\providecommand{\red}{\color{red}}
\providecommand{\enskip}{\quad}
\providecommand{\llbracket}{\lbrack\!\lbrack}
\providecommand{\rrbracket}{\rbrack\!\rbrack}
\providecommand{\euro}{{\text{€}}}
\providecommand{\no}{{\text{n°}}}
\providecommand{\interval}[2]{{[#1;#2]}}
\providecommand{\intervFO}[2]{{[#1;#2[}}
\providecommand{\dotfill}{\ldots}

\providecommand{\ds}{\displaystyle}
\providecommand{\textding}[1]{#1}
\providecommand{\Oij}{\(\mathrm{O};\vec{\imath},\vec{\jmath}\)}
\providecommand{\Oijk}{\(\mathrm{O};\vec{\imath},\vec{\jmath},\vec{k}\)}
\providecommand{\Ouv}{\(\mathrm{O};\vec{u},\vec{v}\)}
\providecommand{\intervFF}[2]{\ensuremath{[#1;#2]}}
\renewcommand{\euro}{\ifmmode\text{\texteuro}\else\texteuro\fi}
\DeclareUnicodeCharacter{00B0}{\ifmmode^{\circ}\else\textdegree\fi}
\DeclareUnicodeCharacter{20AC}{\ifmmode\text{\texteuro}\else\texteuro\fi}
% Énoncés : Éducation nationale. Transcription et corrigés : APMEP, Denis Vergès et contributeurs. Découpage et classement : CoopMaths, Sébastien Lozano, Éric Elter et contributeurs. Adaptation de la présentation : Maths Explorer.
% CC BY-SA : https://coopmaths.fr/www/about/
% https://coopmaths.fr/alea/static/e3c/2020/tex/e3c_2020_00_sujet57_4_cor.tex
\begin{document}
\noindent\textit{APMEP / CoopMaths. CC BY-SA. Présentation adaptée par Maths Explorer.}\par\medskip
\subsection*{1.}

On a :
\[
AB^2 = (1 - (-2))^2 + (2 - 1)^2 = 9 + 1 = 10.
\]
Le carré du rayon est égal à 10.

\[
M(x\,;\,y) \in \mathcal{C} \iff AM^2 = 10 \iff (x + 2)^2 + (y - 1)^2 = 10.
\]

\subsection*{2.}

Avec \(\overrightarrow{AB} \begin{pmatrix} 3 \\ 1 \end{pmatrix}\) et \(\overrightarrow{AE} \begin{pmatrix} 2 \\ -6 \end{pmatrix}\), on a :
\(\overrightarrow{AB} \cdot \overrightarrow{AE} = 6 - 6 = 0\).

\subsection*{3.}

La question précédente montre que les vecteurs \(\overrightarrow{AB}\) et \(\overrightarrow{AE}\) sont orthogonaux, donc les droites \((AB)\) et \((AE)\) sont perpendiculaires.

\subsection*{4.}

On a donc :
\begin{align*}
&M(x\,;\,y) \in (AE) \\
\iff &\overrightarrow{AM} \cdot \overrightarrow{AB} = 0 \\
\iff &3(x + 2) + 1(y - 1) = 0 \\
\iff &3x + y + 5 = 0.
\end{align*}

\subsection*{5.}

Un point \(M(x\,;\,y)\) est commun à la droite \((AE)\) et au cercle \(\mathcal{C}\) si ses coordonnées vérifient leurs équations et donc le système :
\begin{align*}
&\begin{cases}
(x + 2)^2 + (y - 1)^2 = 10 \\
3x + y + 5 = 0
\end{cases} \\
\iff &\begin{cases}
(x + 2)^2 + (-3x - 5 - 1)^2 = 10 \\
y = -3x - 5
\end{cases} \\
\iff &\begin{cases}
(x + 2)^2 + (-3x - 6)^2 = 10 \quad (1)\\
y = -3x - 5
\end{cases}
\end{align*}
\((1)\) donne :
\begin{align*}
x^2 + 4x + 4 + 9x^2 + 36x + 36 &= 10 \\
10x^2 + 40x + 30 &= 0 \\
x^2 + 4x + 3 &= 0.
\end{align*}

L'équation \(x^2 + 4x + 3 = 0\) a une racine évidente : \(-1\) et comme le produit des racines est égal à 3, l'autre racine est \(-3\).

En remplaçant successivement \(x\) par \(-1\) puis par \(-3\) dans l'équation \(y = -3x - 5\), on obtient deux points communs au cercle et à la droite : les points \((-1 \,;\, -2)\) et \((-3 \,;\, 4)\).

\end{document}
